\documentclass[output=paper,colorlinks,citecolor=brown]{langscibook}
\ChapterDOI{10.5281/zenodo.18845598}
\author{Rainer Hofmann\orcid{}\affiliation{} and Moritz Malkmus\orcid{}\affiliation{} and Leonie Östermann\orcid{}\affiliation{}}
\title{Heritage language education. A brief legal overview}
\abstract{This article provides an overview of the legal framework for heritage language education in the federal state of Hesse, with a special focus on the particularities that may arise with regard to accommodating the linguistic needs of persons with a migratory background, especially of those who wish to learn a non-official language or a language that is considered to be a minority language in their country of origin. The analysis indicates that current legislation does not adequately address the situation of these students. Drawing on good practices identified, e.g. by the Council of Europe's Advisory Committee on the  Framework Convention for the Protection of National Minorities, it will be discussed which aspects could be considered (\textit{de lege ferenda}) to design such educational programs in a meaningful way. While not questioning the central importance of acquiring a sound knowledge of the language of instruction, this article argues that students who self-identify as members of certain minority groups and who wish to learn their heritage language, should, under certain circumstances, be given the opportunity to develop this part of their distinct identities, regardless of the formal status of that language.}

\IfFileExists{../localcommands.tex}{
   \addbibresource{../localbibliography.bib}
   \usepackage{tabularx,multicol}
\usepackage{url}
\urlstyle{same}
\usepackage{textcomp}
 
\usepackage{langsci-optional}
\usepackage{langsci-lgr}
\usepackage{langsci-gb4e}
\usepackage{langsci-branding}

% \usepackage{fontspec}
\babelprovide{georgian} % Georgian language
\babelprovide{armenian} % Armenian language
\babelfont[georgian,armenian]{rm}{DejaVu Sans}
\newfontfamily\georgianfont{DejaVu Sans}
\newfontfamily\armenianfont{DejaVu Sans}

\usepackage{dialogue}
\usepackage{attrib}
\usepackage{attrib} % speaker labels in Koca's chapter
\usepackage{phoenician} % alf, bet in Mengozzi's chapter

   

\makeatletter
\let\thetitle\@title
\let\theauthor\@author
\makeatother

\newcommand{\togglepaper}[1][0]{
   \bibliography{../localbibliography}
   \papernote{\scriptsize\normalfont
     \theauthor.
     \titleTemp.
In Saloumeh Gholami \& Ulrich Mehlem (eds.), \textit{Identity formation, language and migration dynamics: Minorities in the MENA region and the German diaspora}, Berlin: Language Science Press [preliminary page numbering]
   }
   \pagenumbering{roman}
   \setcounter{chapter}{#1}
   \addtocounter{chapter}{-1}
}

\newbool{bookcompile}
\booltrue{bookcompile}
\newcommand{\bookorchapter}[2]{\ifbool{bookcompile}{#1}{#2}}


\newcommand{\georgian}[1]{{\georgianfont #1}}
\newcommand{\armenian}[1]{{\armenianfont #1}}

\newcommand{\textarab}[1]{{\arabicfont \beginR #1} \endR}
\newcommand{\texthebrew}[1]{{\hebrewfont \beginR #1} \endR}
\newcommand{\textsyriac}[1]{{\syriacfont \beginR #1} \endR}


\newfontfamily\phoenicianfont{Tuffy Regular.ttf}

   %% hyphenation points for line breaks
%% Normally, automatic hyphenation in LaTeX is very good
%% If a word is mis-hyphenated, add it to this file
%%
%% add information to TeX file before \begin{document} with:
%% %% hyphenation points for line breaks
%% Normally, automatic hyphenation in LaTeX is very good
%% If a word is mis-hyphenated, add it to this file
%%
%% add information to TeX file before \begin{document} with:
%% %% hyphenation points for line breaks
%% Normally, automatic hyphenation in LaTeX is very good
%% If a word is mis-hyphenated, add it to this file
%%
%% add information to TeX file before \begin{document} with:
%% \include{localhyphenation}
\hyphenation{
    par-a-digm
    Ver-to-vec
    Go-go-lin
    sa-gen
    mo-misch
    ant-wor-te-te
    dix-wa-zim
    spre-chen
    B-Verf-GE
    Swe-dish
    re-mains
    schrei-ben
    di-pey-vin
    Chris-ten-sen
    Wes-sen-dorf
    Ja-va-khi-shvi-li
    ana-lysis
}

\hyphenation{
    par-a-digm
    Ver-to-vec
    Go-go-lin
    sa-gen
    mo-misch
    ant-wor-te-te
    dix-wa-zim
    spre-chen
    B-Verf-GE
    Swe-dish
    re-mains
    schrei-ben
    di-pey-vin
    Chris-ten-sen
    Wes-sen-dorf
    Ja-va-khi-shvi-li
    ana-lysis
}

\hyphenation{
    par-a-digm
    Ver-to-vec
    Go-go-lin
    sa-gen
    mo-misch
    ant-wor-te-te
    dix-wa-zim
    spre-chen
    B-Verf-GE
    Swe-dish
    re-mains
    schrei-ben
    di-pey-vin
    Chris-ten-sen
    Wes-sen-dorf
    Ja-va-khi-shvi-li
    ana-lysis
}

   \boolfalse{bookcompile}
   \togglepaper[4]%%chapternumber
}{}
\usepackage{textcomp}
\begin{document}
\maketitle

\section{Preliminary remarks}\label{CH3.ss.1}

The notion of heritage language education -- as is commonly used in other disciplines \citep[3]{trifonas_aravossitas2018} -- does not amount to a uniform concept in law. In a legal context, reference is often made to the teaching of ``mother tongues'',\footnote{%
Cf. Article 3 of Council Directive 77/486/EEC of 25 July 1977 on the education of the children of migrant workers (OJ L 199, 6.8.1977, p. 32--33, \citealt{eec1977}); \textit{Modersmålsundervisning}, Chapter 10, Section 7 of the Swedish \isi{Education} Act (\textit{Skollag}, SFS 2010:800, \citealt{skollag2010}).} heritage languages,\footnote{%
\textit{Unterricht in der Herkunftssprache}, cf. Section 8a para. 3, sentence 1 of the Hessian School Act (\textit{Hessisches Schulgesetz}) in its version of 1 August 1997 (expired on 31 July 1999), GVBl. 1997 I S. 143, \citep{gvbl1997}.} languages and cultures of origin,\footnote{%
Cf. the \ili{French} Enseignement Langue et Culture d’origine (ELCO), a program, which has been organized since the 1970s.} foreign languages,\footnote{%
Cf. the \ili{French} Enseignements Internationaux de Langues étrangères (EILE), the successor program, which has been successively replacing ELCO since 2016.} or \isi{minority} languages.\footnote{%
Cf. Article 14 of the Council of Europe’s Framework Convention for the Protection of \isi{National} Minorities (FCNM, ETS 157, \citealt{fcnm1995}); Article 8 of the European Charter for Regional or Minority Languages (ECRML, ETS 148, \citealt{ecrml2024}).}
% \todo[inline]{Check insertion of publication citation after legal citation in footnotes 1, 2, 5 and throughout.
% Check insertion of words \textit{Modersmålsundervisning} (fn. 1) and 
% \textit{Unterricht in der Herkunftssprache} (fn. 2).
% Delete cf. in footnotes 1--5 and maybe elsewhere.}
While the above terminology encompasses a variety of different educational programs, the following reflections will start by approaching this issue from a linguistic perspective (\citealt[23]{mehlhorn2020}, \citealt[8]{polinsky2015}) that places the term \textit{heritage language} and its acquisition\is{Language!acquisition} in relation to its speakers. In this sense, ``heritage speakers'' can be defined as 
\largerpage[2]

\begin{quote}
    individuals who were raised in homes where a language other than the dominant community language was spoken and thus possess some degree of \isi{bilingualism} in the heritage language and the dominant language \citep[8]{polinsky2015}.
\end{quote}

This diversity, which exists to varying degrees in all European societies (cf. \citealt{gogolin2002}, \citealt{studer_werlen2012}), requires the linguistic needs of speakers of non-dominant languages to be adequately addressed in the social environment of the dominant community language. In the field of education, where, according to \citet[36]{oecd2021}, around 12\% of students speak a language other than the language of assessment in their home environment, the question arises as to the status to be accorded to these (non-dominant) heritage languages and, in particular, the extent to which states should encourage the acquisition\is{Language!acquisition} of such languages, alongside a sound knowledge of the language of instruction, which is commonly considered a \textit{conditio sine qua non} for educational success. Providing appropriate answers to these questions seems all the more important given the special relationship between language acquisition\is{Language!acquisition} and a person's identity, as illustrated, as illustrated by \citet{edwards2009}, \citet{krumm2020}, and throughout this volume. 

When it comes to the organization of heritage language education, three basic models can be distinguished in Europe: The first model, sometimes referred to as the Scandinavian model, is characterized by heritage language programs that are implemented within the personal and financial capacities of the country of residence; a second cooperation-based model relies on bilateral agreements on teaching in the respective official languages, including the financing and secondment of teachers by the partner state or some form of consular teaching; and a third communitarian model leaves the responsibility for teaching heritage languages to the respective communities (cf. \citealt[82]{reich2017}, \citealt[25]{mehlhorn2020}, \citealt[207--211]{woerfel_kuppers_schroeder2020}, \citealt[5--12]{hofmann_malkmus2022a}). In the federal state of \isi{Hesse}, educational policies\is{Language!policies} have led to a mixture of the first and second option, i.e., state-organized and cooperation-based approaches, with the first option gradually being phased out and responsibility for heritage language teaching being increasingly placed in the hands of partner states and their consulates (cf. \citealt{koller2012}, \citealt[15--18]{hofmann_malkmus2022a}). This seems problematic, not only because of the need to ensure the quality and neutrality of education, but in particular with respect to the situation of non-official languages.\footnote{%
This notion refers to languages that have not acquired official status, even in the states where the corresponding groups habitually reside within the larger population.}

Against this background, the legal situation of pupils who wish to learn a non-official language or a language considered a \isi{minority} language in their country of origin, such as \ili{Kurdish} or \ili{Berber}, will be discussed below. In particular, the question of whether states are obliged to provide corresponding educational programs will be addressed. This will be followed by some reflections on best practices in \isi{minority} language education identified by the Council of Europe's Advisory Committee on the Framework Convention for the Protection of \isi{National} Minorities (ACFC) before concluding with some remarks \textit{de lege ferenda}. Although there is some substantial work on related topics, e.g., by \citet{langenfeld1998}, \citet{schmahl2004}, \citet{medda-windischer2008}, \citet{gogolin_oeter2011}, \citet{guliyeva2013}, \citet{palermo2019}, \citet{jacob-owens2022}, \citet{bayat2023}, the above peculiarities still seem to be an under-researched area of law.

\section{Legal considerations}\label{CH3.ss.2}

\subsection{Domestic law}\label{CH3.ss.2.1}

In principle, the organization of the public school system in Germany falls within the competencies of the federal states (\textit{Länder}) (Article 70 para. 1 of the Basic Law). In the federal state of \isi{Hesse}, the relevant constitutional provisions for the state school system are set out in Articles 55--61 of the Hessian Constitution. Although there are provisions on compulsory education, educational objectives, religious and ideological denominations, and the teaching of history, the Hessian Constitution remains silent on language teaching in public schools. Consequently, in principle, the matter is left to the discretion of the State Parliament (\textit{Landtag}) (cf. Article 56 para. 7 of the Hessian Constitution).

In August 1997, the Hessian \textit{Landtag} introduced Section 8a para. 3 to the Hessian School Act (\textit{Hessisches Schulgesetz}) (GVBl. 1997 I S. 143, \citealt{gvbl1997}), which stated: 

\begin{quote}
    Die Förderung der natürlichen Mehrsprachigkeit erfolgt in der ersten Phase durch Unterricht in der Herkunftssprache, der dem Bedarf und den personellen und sächlichen Möglichkeiten der Schulen entsprechend eingerichtet werden soll. Der Unterricht ist so zu gestalten, daß er zur interkulturellen Bildung und Erziehung beiträgt. In der zweiten Phase kann dieser Unterricht als Wahlpflichtangebot fortgeführt oder die Herkunftssprache als zweite Fremdsprache angeboten werden.\\ 
    \char`[The promotion of natural multilingualism takes place in the first phase \char`[primary education\char`]\ through teaching in the heritage language, which should be organized according to the demand and the personnel and material resources of the schools. The lessons are to be designed in such a way that they contribute to intercultural education and upbringing. In the second phase \char`[of school education\char`], these lessons can be continued as a compulsory elective, or the language of origin can be offered as a second foreign language.\char`]\ (Translation by the authors)
\end{quote}


However, only two years later, in July 1999, this provision expired and was replaced by the following transitional provision (Article 3 para. 1 No. 4 of the Act 30 June 1999 (GVBl. 1999 I S. 354, \citealt{gvbl1999}):

\begin{quote}
    Unterricht in einer Herkunftssprache nach dem bisherigen § 8a Abs. 3 des Schulgesetzes kann im Rahmen bestehender personeller Möglichkeiten fortgeführt werden.Dabei sind Schülerinnen und Schüler, die bei In-Kraft-Treten dieses Gesetzes an dem Unterricht teilnehmen, vorrangig zu berücksichtigen. \\
    \char`[Teaching in a heritage language in accordance with the previous Section 8a para. 3 of the School Act may be continued within the framework of existing personnel capabilities. Priority shall be given to pupils who are attending classes at the time of the entry into force of this Act. (Translation by the authors)\char`]
\end{quote}

This situation led to the gradual transfer of state-organized heritage language education in Hesse to the responsibility of the countries of origin within the framework of the instruction organized by the consulates. Even if the existing state-organized offers may be continued (\textit{kann} […] \textit{fortgesetzt werden}) on the basis of the above-mentioned transitional provisions, the state responsibility for heritage language teaching became a discontinued model (cf. \citealt{koller2012}).

Against this backdrop, the Wiesbaden Administrative Court ruled in 2018 that students do not have a subjective right to the establishment of additional capacities for teaching their heritage language (VG Wiesbaden, Urteil vom 13.11.2018, 6 K 1560/18.WI, \citealt{vg_wiesbaden2019}). In the case at hand, the applicant requested to receive heritage language classes at her elementary school in \textit{Kurmanji},  which the Court considered to be a \ili{Kurdish} \isi{dialect}t. 

While the relevant aspects of EU and \isi{international law} will be addressed below, two observations and one additional remark can already be made with regard to the domestic legal situation: First, the decision of the Administrative Court illustrates that the current Hessian school legislation, as set out in the Hessian School Act and the Regulation on the Timetable for Primary and Lower Secondary Schools (ABl. 2011 p. 653, \citealt{abl2011}), does not meet the requirements for establishing a subjective right of pupils to be provided with a specific heritage language course. This is mainly due to the legal design of the relevant provisions, which expressed the legislator's intention not to give rise to any subjective (enforceable) claims and to transform state-organized heritage language education into a discontinued model.

Second, as the Court further states, such a claim cannot be derived from the constitutional principle of equal treatment (Article 1 of the Hessian Constitution; Article 3 of the Basic Law). In its dimension as a right to equal access to school education, the mentioned principle grants a derivative right to participate in existing programs in a non-arbitrary and non-discriminatory manner (\textit{derivatives Teilhaberecht}). Inequalities may be justified -- to a certain extent -- by reasons of school capacities. In this sense, the Administrative Court considered the different treatment of, for example, \ili{Polish} heritage language speakers to be justified. Furthermore, with reference to the Federal Constitutional Court's (FCC, \textit{Bundesverfassungsgericht}) established case law (BVerfGE 33, 303, \citealt{bverfge1972}), the Administrative Court pointed out that an inherent positive right to participation (\textit{originäres Teilhaberecht}) in education is limited to the existing capacities and is therefore subject to the proviso of the possible (\textit{Vorbehalt des Möglichen}). According to the Court, this entails that the legislator (and to some extent the school administration) carries the primary responsibility to decide on the allocation of available resources. Hence, the principle of equal treatment does not support any claim to design state education in a particular way.

Following the decision of the First Senate of the FCC of 19 November 2021 on the so-called ``Federal pandemic emergency brake II'' (\textit{Bundesnotbremse II}, BVerfGE 159, 355, \citealt{bverfge2021}), the Administrative Court's interpretation of constitutional equality requirements in education can be considered broadly confirmed. Although the FCC addressed a different set of facts, namely the impact of anti-pan\-demic measures on school education (cf. \citealt{donath2023}), the legal dimensions of the right to education, set out by the FCC, are essentially congruent with the basic line of reasoning adopted by the Wiesbaden Administrative Court.

\textit{Third}, the question arises as to whether the domestic legal situation should be assessed differently in light of the 2018 amendments to the Hessian Constitution. According to the newly introduced Article 4 para. 2, sentence 1 of the Hessian Constitution, every child has a ``right to protection and to the promotion of his or her development into an independent and socially competent personality''.\footnote{%
See Article 4 para. 2 of the Hessian Constitution as amended by the Act of 12 December 2018 (GVBl. S. 752, \citealt{gvbl2018}): \textit{Jedes Kind hat das Recht auf Schutz sowie auf Förderung seiner Entwicklung zu einer eigenverantwortlichen und gemeinschaftsfähigen Persönlichkeit.} It should be noted, however, that the Wiesbaden Administrative Court has already dealt with questions of the best interests of the child and nevertheless concluded that, from this point of view, no entitlement to the requested language instruction can be recognized, unless there are unreasonable circumstances which seriously impair the best interests of the child.} Given the importance of language acquisition\is{Language!acquisition} for the development of the child's own identity, it could be argued that the right to promotion of the development of a child's personality includes, to some extent, learning her/his heritage language, which should be facilitated by state education. Since, according to the legislative documents (Drucks. 19/5710 p. 2, \citealt{hesse2017}), the \textit{Landtag} was guided by the Convention on the Rights of the Child (CRC, 1577 UNTS 3, \citealt{un_crc_1989}) when incorporating this fundamental right of children into the Hessian Constitution (cf. \citealt{donath2020}), it seems reasonable to assume that the interpretation of this new constitutional provision is informed, to a large extent, by the principles set out in the CRC. This includes, as will be shown below, respect for the child's cultural identity and language as an objective of the state school system (Article 29 para. 1 lit. c of the CRC). The structure of the Hessian Constitution, in particular the placement of Article 4 para. 2 outside the regulatory context of the state school system (Articles 55--61), does not necessarily contradict this interpretation either. However, in light of the inherent limitations described above, it is doubtful whether a subjective entitlement to the desired language instruction would automatically arise. Even if Article 4 para. 2 of the Hessian Constitution were to be interpreted as linking language acquisition\is{Language!acquisition} with personal development, the desire to learn a particular heritage language could not be granted absolute priority (in the sense of subjective entitlements) due to the proviso of the possible and the primary responsibility of the legislator to organize school education by specific legislation. Nevertheless, it can be argued that the objective-legal dimension of Article 4 para. 2 of the Hessian Constitution obliges state authorities to take due account of the linguistic needs of these children in their discretionary decisions (alongside the other objectives of the state education system and within the limits mentioned above) and, where appropriate, provide such education within the framework of state-organized schooling.

\subsection{EU law}\label{CH3.ss.2.2}

With regard to EU law, it should be recalled that under the fundamental principle of conferral, ``the Union shall act only within the limits of the competences conferred on it by the Member States in the Treaties to attain the objectives set out therein'' (Article 5 para. 2 of the Treaty on European Union, TEU). In the field of education, Article 165 para. 1 of the Treaty on the Functioning of the European Union (TFEU) emphasizes the ``responsibility of the Member States for both the content of teaching and the organisation of their education systems'', while further ``harmonisation of the laws and regulations of the Member States'' by Union acts is excluded under Article 165 para. 4 TFEU. This means that education policies\is{Language!policies} are largely reserved for the Member States (cf. \citealt[220--222]{guliyeva2013}; \citealt[1426 et seq.]{garben2019}). 
Despite this, in 1977 the Council made use of an enabling provision\footnote{%
Article 49 of the Treaty establishing the European Economic Community (now Article 46 TFEU).} on the free movement of workers and adopted Directive 77/486/EEC on the education of the children of migrant workers (OJ L 199, 6 August 1977, p. 32, \citealt{eec1977}). Article 3 of this Directive provides: 

\begin{quote}
    Member States shall, in accordance with their national circumstances and legal systems, and in cooperation with States of origin, take appropriate measures to promote, in coordination with normal education, teaching of the mother tongue and culture of the country of origin for the children referred to in Article 1.\footnote{%
    A similar provision can be found in Article 19 para. 12 of the revised European Social Charter (ETS 163, \citealt{coe1996}) or Article 15 of the European Convention on the Legal Status of Migrant Workers (ETS 93, \citealt{coe1977}). Cf. \citealt[30--31]{hofmann_malkmus2022a}.}
\end{quote}

Despite existing implementation shortcomings, Council  Directive 77/486/EEC remains one of the few normative bases for heritage language teaching in the EU Member States (cf. \citealt{koller2012}, \citealt[221]{guliyeva2013}). However, with regard to the present case, the Administrative Court of Wiesbaden, in its above-mentioned judgment (6 K 1560/18.WI), pointed out that Directive 77/486/EEC only applies to children living in a host Member State who are dependants of a worker who is a national of \textit{another} EU Member State and in addition, that it only covers education in the language of origin of an EU Member State, thus excluding immigrants from non-EU countries, such as \isi{Morocco}, Syria, Türkiye, and several non-official languages. This underlines the limited scope of the Directive (cf. \citealt{koller2012}).

Beyond this, and despite various references to ``\isi{minorities}'' (Article 2 TEU),\footnote{%
See also Article 21 para. 1 of the EU Charter of Fundamental Rights \citep{cfr2012}.} or ``linguistic diversity'' (Article 22 of the EU Charter of Fundamental Rights, \citealt{cfr2012}), \isi{EU law} does not provide for any explicit basis to take further legal action in the field of \isi{minority rights} or \isi{minority} languages (\citealt[50]{toggenburg2018}, \citealt{hofmann_malkmus2022b}). Notwithstanding the EU's particular interest in \isi{minority} issues in the context of the accession process \citep[304]{malkmus2023b}, the Union remains reluctant about its internal \isi{minority} policy, as recently demonstrated by the setback of the ``Minority SafePack'' initiative \citep{malkmus2024}. It is, therefore, unlikely that the legal situation described above will change in the near future.

\subsection{International law}\label{CH3.ss.2.3}

In \isi{international law}, following \citet[38 et seq.]{henrard2013}, for example, a distinction can be made between \isi{minority}-specific and non-\isi{minority}-specific provisions in order to narrow down the question under consideration here. The first group of norms includes, \textit{inter alia}, Article 27 of the International Covenant on Civil and Political Rights (ICCPR, 999 UNTS 171, \citealt{un_iccpr_1966}), Article 14 para. 2 of the Framework Convention for the Protection of \isi{National} Minorities (FCNM, ETS 157, \citealt{fcnm1995}) and Article 8 of the European Charter for Regional or Minority Languages (ECRML, ETS 148, \citealt{ecrml2024}). In the present context, the situation of persons with a migratory background, particularly those seeking to learn a non-official language, raises the question of whether they fall within the personal scope of the above provisions at all. From a textual analysis, it is apparent that only the ECRML does not apply \textit{a priori} to the ``languages of migrants'' (Article 1 lit. a), while the two above instruments are somewhat ambiguous in this respect. Although there is no uniform definition of the concept of \isi{minorities} in \isi{international law}, it is widely assumed in state practice that these persons do not fall within the scope of the relevant treaties, while international expert bodies have adopted more nuanced approaches, seeking to include these groups, where appropriate, on a situation-specific basis (cf. \citealt[para. 4--5]{hofmann2007_minorities}, \citealt{malkmus2023a}, \citealt{hofmann_malkmus2023a}). In particular, the ACFC is considering, on an article-by-article basis, whether not only so-called ``old'' or autochthonous \isi{minorities} but also ``new'' or allochthonous \isi{minorities} (i.e. persons with a migratory background) could benefit from the FCNM \citep{palermo2019}. However, as explained in the ACFC's 4th Thematic Commentary (ACFC/56DOC(2016)001, para. 79, 83, \citealt{acfc2016}), the wording of Article 14 of the FCNM contains an inherent limitation of the right to education or instruction in the \isi{minority} language to ``areas inhabited by persons \isi{belonging} to national \isi{minorities} traditionally or in substantial numbers'' (Article 14 para. 2 FCNM, cf. \citealt{busch2018}: 260--262, \citealt{hofmann2023a}: 35). Nevertheless, the ACFC encourages states parties to promote the enjoyment of this right

\begin{quote}
    in situations where the conditions are not formally met but where implementation would serve to promote an open society, where multilingualism is encouraged as a reflection of diversity (ACFC/{\allowbreak}56DOC{\allowbreak}(2016)001, para. 79, \citealt{acfc2016}; \citealt[24]{palermo2019}).
\end{quote}

A similar approach could already be observed earlier in the UN Working Group on Minorities' commentary to the 1992 UN Minorities Declaration (A/{\allowbreak}RES/{\allowbreak}47/{\allowbreak}135), which, as a non-binding instrument, aims to specify the rights set out in Article 27 ICCPR \citep[875]{hofmann_malkmus2023a}. With regard to the right to mother tongue education, the Working Group commented that there would be ``no reason'' to treat new \isi{minorities} differently from old \isi{minorities} if they ``settle compactly together in a region of the country and in large number'' (E/CN.4/{\allowbreak}Sub.2/{\allowbreak}AC.5/{\allowbreak}2005/2, para. 64, \citealt{wgm2005}). Furthermore, the UN Special Rapporteur on \isi{minority} issues (A/HRC/43/47, para. 53 et seq., \citealt{unsr2020}) also supports an inclusive reading of the educational rights of children \isi{belonging} to non-recognized \isi{minority} groups to learn their mother-tongue.

Against this background, although the right to heritage language education for new \isi{minorities} has not (yet) acquired a clear normatively binding status in international \isi{minority rights} law, it is apparent from the above-mentioned statements that, in certain circumstances, in order to promote linguistic diversity and integration, States may focus less on formal distinctions and consider, where feasible, whether there is a \isi{minority} that could benefit from the right to heritage language education.

The second set of provisions encompasses the right to education\footnote{%
E.g. Article 13 para. 1 of the International Covenant on Economic, Social and Cultural Rights (ICESCR, 993 UNTS 3, \citealt{un_icescr_1976}), Article 2 of the Protocol to the ECHR (ETS 9, \citealt{coe1952}), Article 28 para. 1 in conjunction with Article 29 para. 1 lit. c of the Convention on the Rights of the Child (CRC, 1577 UNTS 3).} and the prohibition of discrimination.\footnote{%
E.g. Article 26 of the ICCPR, Article 5 lit. e of the International Convention on the Elimination of All Forms of Racial \isi{Discrimination} (ICERD, 660 UNTS 195, \citealt{un_icerd_1965}), Article 14 of the European Convention for the Protection of Human Rights and Fundamental Freedoms (ECHR, ETS 5, \citealt{coe1950}) and Article 1 of Protocol No. 12 to the ECHR (ETS 177, \citealt{coe2000}).} Without being able to go into all the details, the following three points should be highlighted: First, under Article 2 of the Protocol to the ECHR, the right to education, as interpreted by the case law of the European Court of Human Rights (e.g. judgment [GC] of 10 May 2001, appl. no. 25781/94, Cyprus v. Türkiye, para. 277, \citealt{echr2010}; judgment of 14 September 2023, appl. no. 56928/19, 7306/20, 11937/20, Valiullina and Others v. Latvia, para. 135, \citealt{echr2023}), does not require school education to be provided in a particular language (cf. \citealt{bitter2022}: para. 14). This reading of the right to education was also adopted by the Wiesbaden Administrative Court (6 K 1560/18).

A more dynamic interpretation could be observed in several opinions of the Committee on Economic, Social and Cultural Rights (CESCR), which, for example, in its observations on Greece (E/C.12/1/Add.97, para. 28, 50, \citealt{cesc2004}), criticized the state authorities for not providing sufficient education in the mother tongue of the \ili{Turkish} and Roma \isi{minorities} \citep[1134]{saul2014}. It is worth mentioning, though, that the groups to which the CESCR refers in its observations are considered to be long-established \isi{minorities} within Greece, leaving open the question of whether this would also apply to new \isi{minorities}. Furthermore, the Committee has been concerned with a situation where a ``high percentage'' of these children were ``not enrolled in school, or drop out at a very early stage of their schooling'' (E/C.12/1/Add.97, para. 28, \citealt{cesc2004}), i.e. where not providing education in such languages could be seen equal to impeding access to education.

On the other hand, the Committee on the Elimination of Racial \isi{Discrimination} (CERD/C/DNK/CO/18-19, para. 16, \citealt{cerd2010}) and later the Committee on Economic, Social and Cultural Rights (E/C.12/DNK/CO/6, para. 68, \citealt{cescr2019}) have been remarkably clear in their criticism of the distinction between children of European origin and other children, as is the case, for example, in the Danish heritage language education \citep[37--38]{hofmann_malkmus2022a}. Therefore, the second observation is that any distinction between official EU languages and other languages must be justified, and purely formalistic reasons stemming from the EU origin appear insufficient. Appropriate reasons for differentiation still appear to be school capacities, availability of accredited teachers and individual demand.

Third, apart from those subjective legal safeguards, States are obliged to take due account of the linguistic needs of children, including the promotion of their heritage language, when organizing State education. This obligation is of an objective nature and arises from Article 29 para. 1 lit. c of the CRC:

\begin{quote}
    States Parties agree that the education of the child shall be directed to: […] The development of respect for the child's parents, his or her own cultural identity, language and values, for the national values of the country in which the child is living, the country from which he or she may originate, and for civilizations different from his or her own […].
\end{quote}

In this context, the Committee on the Rights of the Child (CRC/C/CHN/CO/3-4, para. 74, \citealt{crc2013}; CRC/C/CYP/CO/3-4, para. 44--45, \citealt{crc2012cyprus}; CRC/C/TUR/CO/2-3, para. 58, \citealt{crc2012turkey}) has encouraged states parties in the past (regardless of the formal status of the language) to introduce or expand mother tongue education programs where appropriate (\citealt[1137]{lundy_tobin2019}, \citealt[303]{schmahl2020}). Interestingly, the Committee does not seem to want to make this dependent on the formal status of a \isi{minority} or language \citep[1137]{lundy_tobin2019}, as the concluding observations on Türkiye suggest: 

\begin{quote}
    Unavailability of education in languages other than \ili{Turkish} and languages of recognized \isi{minorities},\footnote{%
    Due to the \ili{Turkish} interpretation of the 1923 Lausanne Peace Treaty (and its annexes), these are the languages of religious \isi{minorities} such as Jews, Greeks and Armenians \citep[65]{scherzberg2014}. See also the reservation of Türkiye under Article 27 ICCPR.} presenting educational disadvantages for children of non-recognized \isi{minorities} whose mother tongue is not \ili{Turkish} (CRC/C/TUR/CO/2-3, para. 58, \citealt{crc2012turkey}).
\end{quote}

\subsection{Conclusion}\label{CH3.ss.2.4}

It can be concluded that, beyond the right to non-discriminatory access to existing educational programs, there is currently no inherent positive right to education in non-official heritage languages, unless explicitly granted by specific legislation. However, states are required to take due account of the linguistic needs of children who wish to learn such languages as part of their distinct identities when designing language education. Even if this does not give rise to subjective (enforceable) rights, there is an objective obligation to consider whether this desire can be accommodated within existing capacities. This includes, where appropriate, facilitating the teaching of heritage languages.

\section{Good practices}\label{CH3.ss.3}

Following this interpretation, the present section outlines some good practices from which states could draw inspiration to fulfil the above obligation. The Swedish legislation on heritage language education, as laid down in the \isi{Education} Act (\textit{Skollag}, SFS 2010:800, \citealt{skollag2010}) and the Language Act (\textit{Språklag}, SFS 2009:600, \citealt{spraklag2009}), may serve as an example (cf. ACFC/OP/IV(2017)004, para. 91, \citealt{acfc2017}; \citealt{salo2018}; \citealt{hofmann_malkmus2022a}: 5). According to Section 14 of the Language Act, every person whose mother tongue is not one of the national, \isi{minority} or sign languages shall ``be given the opportunity to develop and use [his or her] mother tongue''. Mother tongue teaching (\textit{Modersmålsundervisning}) is therefore part of the Swedish curriculum for the different types of school. For primary schools, for example, this is stipulated in Chapter 10, Section 7 of the \isi{Education} Act:

\begin{quote}
    A pupil who has a guardian with a mother tongue other than \ili{Swedish} shall be offered mother tongue teaching in this language if
    \begin{enumerate}
        \item the language is the pupil's daily language of interaction at home, and
        \item the pupil has basic knowledge of the language.
        \item[] [...]
    \end{enumerate}
    The Government or the authority designated by the Government may issue regulations on mother tongue teaching. Such regulations may mean that mother tongue teaching shall be offered in a language only if a certain number of pupils want such instruction in that language. [...] (Translation by the authors)
\end{quote}

The details of mother tongue teaching are laid down in Chapter 5 of the School Regulation (\textit{Skolförordning}, SFS 2011:185, \citealt{skolforordning2011}) of the Ministry of \isi{Education}: 

\begin{quote}
    8 § Mother tongue teaching may be organized
    \begin{enumerate}
        \item as a language choice in compulsory school and special school,
        \item as the pupil's choice,
        \item within the framework of the school's choice, or
        \item outside the guaranteed teaching time.
        \item[] […]
    \end{enumerate}
    10 § A principal is obliged to organize mother tongue teaching in one language only if
    \begin{enumerate}
        \item at least five pupils who are to be offered mother tongue teaching in the language wish to receive such teaching, and
        \item there is a suitable teacher.
        \item[] […] (Translation by the authors)
    \end{enumerate}
\end{quote}

It should be noted that the Swedish legislation has chosen a dual structure, with some basic legal provisions in the Language and \isi{Education} Acts, i.e. formal legislation adopted in the \textit{Riksdagen}, and the details further specified by a ministerial regulation, which leaves these details open to change (outside parliamentary procedures). This legislative approach reflects a compromise between the linguistic needs of the children and certain concessions for flexibility and practical reasons, such as finding suitable teachers or reaching the required number of pupils, which can still be major obstacles, especially for less widely spoken languages. This situation certainly leaves plenty of room for debates about the appropriate parameters, also in view of the considerable financial demands. Still, the fact that the teaching of heritage languages is enshrined in a parliamentary act that focuses on those mother tongues that are actually being used (rather than on their legal status) is an added value in itself, even if the details are left to the Ministry of \isi{Education} \citep[5--7, 12]{hofmann_malkmus2022a}.

Another source of inspiration might be the experience of the ACFC, bearing in mind, of course, that there are structural and legal differences between the teaching of heritage languages on the one hand and \isi{minority} languages on the other. Some models established in a \isi{minority} context, such as the trilingual school system in two \ili{Ladin} valleys of South Tyrol, which integrates \ili{Italian}, \ili{German}, and \ili{Ladin} in the classroom, have proven particularly successful in specific geographical settings where the population consists of distinct \isi{minority} groups. These models do not seem to be easily transferable to the situation of children who speak various less widespread or migrant languages, as it seems difficult to integrate a large number of additional languages into these kind of models  (cf. \citealt{carla2018}: 1109, \citealt{vettori2021}: 241 et seq.). For similar reasons, models employing the \isi{minority} language as the medium of instruction, as identified within the Council of Europe \citep{prina_dunbar_gurbo2020}, are unlikely to be effective for the purposes outlined above. However, the practice of the ACFC on the teaching of \isi{minority} languages has provided some valuable lessons that may be worth considering here. The first could be seen in the ACFC's proactive approach, which includes making people aware of the educational opportunities available to them, as only such awareness can create demand (ACFC/25DOC[2006]002, p. 26, \citealt{acfc2006}; ACFC/44DOC[2012]001rev, para. 71, \citealt{acfc2012}). Second, state authorities should create explicit legal bases for (\isi{minority}) language education, which also pay special attention to the situation of languages of numerically smaller groups (ACFC/44DOC[2012]001rev, para. 70, \citealt{acfc2012}). Third, there ``should be no mutually exclusive choice between the learning of a \isi{minority} language or the official language(s)'' (ACFC/44DOC[2012]001rev, para. 72, \citealt{acfc2012}). In this sense, it is important not to make the acquisition\is{Language!acquisition} of the heritage language dependent on a certain level of knowledge of the official language. Fourth, \isi{minority} language education should be included into the mandatory curriculum of public schools and numerical thresholds for the establishment of such classes should be set moderately (ACFC/44DOC[2012]001rev, para. 73--74, \citealt{acfc2012}). Finally, and this cannot be overemphasized, there is a significant need for teacher training and teaching materials (ACFC/44DOC[2012]001rev, para. 76, \citealt{acfc2012}). In addition, the situation of particularly small \isi{minority} groups requires special treatment, which may make it necessary either to resort to cooperation-based solutions or to utilize the advantages of digital teaching media in order to create appropriate offers for language teaching \citep[40]{hofmann2023a}. In particular, where the state cannot reasonably be expected to provide classroom-based education -- for example, due to low demand or lack of staff -- the extent to which such gaps can at least be filled by digital learning opportunities could be further explored (ideally in coordination with the relevant communities). In this respect, the objective duty described above can be understood as a best effort obligation.

\largerpage
Further impetus could be provided by the concept of ``plurilingual and intercultural education'', which is promoted notably by the Committee of Ministers of the Council of Europe (Recommendation CM/Rec[2022]1, \citealt{cm2022}) or the Parliamentary Assembly's Recommendation 1740 (2006) on the place of mother tongue in school education \citep{pa2006}.

\section{Concluding remarks \textit{de lege ferenda}}\label{CH3.ss.4}

Although there is currently no inherent positive right to education in non-of\-fi\-cial or immigrant (heritage) languages, unless explicitly granted by specific legislation, states are required to take due account of the linguistic needs of children who wish to learn such languages as part of their distinct identities when designing language education. Where resources are limited, it may be necessary to impose certain constraints. These should not be made unilaterally at the expense of non-official languages. Instead, programs that are to be designed could take the following points into account:

\begin{enumerate}
    \item The local population of speakers of a particular heritage language and the resulting demand; and
    \item the availability of suitable teachers and teaching materials.
    \item Authorities should pursue a proactive approach in order to facilitate access to existing programs.
    \item The scarcity of resources may make it necessary to explore new ways of providing an appropriate educational experience for the pupils concerned, for example, by creating networked programs or exploring the benefits of digital media.
\end{enumerate}
   
The framework for such educational programs could, as is the case in \isi{Sweden}, be laid down in a corresponding legal basis in the Hessian School Act, which could delegate further details to a regulation addressing the considerations outlined above and allowing for necessary adjustments. In any case, the contributions in this volume have illustrated the need to take due account of the special significance of language acquisition\is{Language!acquisition} for the development of one's own identity in the school context. Teaching heritage languages in some of the forms described above could make a significant contribution in this respect.

\section*{Abbreviations}
\small
\begin{tabularx}{.49\textwidth}{lQ}
ABl. & Amtsblatt \\
ACFC & Advisory Committee on Framework Convention for the Protection of \isi{National} Minorities \\
Add. & Addendum \\
appl. no. & application number \\
BVerfGE & Entscheidungen des Bundesverfassungsgerichts \\
CERD & Committee on the Elimination of Racial \isi{Discrimination} \\
CESCR & Committee on Economic, Social and Cultural Rights \\
CM & Committee of Ministers \\
CRC & Convention on the Rights of the Child \citep{un_crc_1989}\\
Drucks. & Drucksache \\
ECHR & European Convention on Human Rights \citep{coe1952}\\
ECRML & European Charter for Regional or Minority Languages \citep{ecrml2024}\\
EEC & European Economic Community \\
EILE & Enseignement Internationaux de Langues Étrangère \\
ELCO & Enseignement Langue et Culture d’origine \\
et seq. & \textit{et sequentia} `and the following' \\
EU & European Union \\
FCC & Federal Constitutional Court \\
FCNM & Framework Convention for the Protection of \isi{National} Minorities \citep{fcnm1995}\\
\end{tabularx}%
\begin{tabularx}{.49\textwidth}{lQ}
GC & Grand Chamber of the European Court of Human Rights \\
GVBl. & Gesetz- und Verordnungsblatt \\
HRC & United Nations Human Rights Committee \\
ICCPR & International Covenant on Civil and Political Rights \citep{un_iccpr_1966}\\
ICERD & International Convention on the Elimination of All Forms of Racial \isi{Discrimination} \citep{un_icerd_1965}\\
ICESCR & International Covenant on Economic, Social and Cultural Rights \citep{un_icescr_1976}\\
lit. & \textit{litera} `letter' \\
No. & Number \\
OECD & Organisation for Economic Co-operation and Development \\
OJ & Official Journal \\
p. & page \\
para. & paragraph \\
Rec & Recommendation \\
RES & Resolution \\
SFS & Svensk författningssamling \\
TEU & Treaty on European Union \citep{teu1992}\\
TFEU & Treaty on the Functioning of the European Union \citep{tfeu2016}\\
UN & United Nations \\
UNTS & United Nations Treaty Series \\
VG & Verwaltungsgericht \\
ZP & Zusatzprotokoll \\
  & \\
  & \\
  & \\
  & \\
\end{tabularx}

\normalsize
\sloppy
\printbibliography[heading=subbibliography,notkeyword=this]
\end{document}
